\section{L'éducation et les pratiques autochtones}\label{leducation-et-les-pratiques-autochtones}

\stanceinfo{\textbf{FCÉG CCLI 2024} [2024-01-05]}{2024-01-05}

\officialstance{La Fédération canadienne étudiante de génie reconnaît les torts causés aux peuples autochtones et admet qu'une véritable réconciliation est un processus continu. Le rôle de l'ingénieur dans la réconciliation comprend la reconstruction de notre relation avec les peuples autochtones. Pour intégrer la perspective autochtone dans la résolution des problèmes d'ingénierie, les programmes d'études en ingénierie doivent inclure le savoir autochtone, et les institutions doivent déployer des efforts supplémentaires pour accroître la représentation autochtone dans le domaine de l'ingénierie.}

\stanceheading{L'enjeu}

Les questions autochtones, les connaissances et la réconciliation sociale devraient être prioritaires dans la formation des ingénieurs. Les peuples autochtones ont subi des épreuves tout au long de l'histoire du Canada, notamment la dépossession de leurs terres ancestrales et le fait de se voir refuser la possibilité de jouer le rôle de consultant principal pour les questions liées à leur patrimoine. Une perspective intersectionnelle est fondamentale pour la compréhension et le respect mutuel qui doivent être établis afin de progresser vers un changement positif.

La décolonisation s'efforce de créer un espace sûr pour l'ensemble du personnel, du corps enseignant et des étudiants autochtones dans les institutions; plus important encore, elle protège les populations autochtones locales afin de promouvoir la collaboration et les enseignements locaux. En raison du manque d'intégration autochtone locale, il existe une disparité entre l'enseignement de l'histoire et le programme d'ingénierie traditionnel. L'éducation autochtone localisée enseigne comment communiquer efficacement et respectueusement avec les parties prenantes autochtones et indique quand s'engager et quelles questions doivent être abordées en fonction de la région géographique.

Au sein de la FCEG, nous pensons que l'éducation et l'ingénierie autochtone devraient mettre l'accent sur les interactions avec les peuples autochtones. Nous pensons que le BCAPG ne dispose pas actuellement d'une formation autochtone localisée en tant qu'attribut des diplômés, afin d'observer une perspective intersectionnelle et de créer des ingénieurs bien équilibrés, capables de traduire les connaissances autochtones et de communiquer avec les communautés indigènes dans le cadre de leur profession d'ingénieur. Des changements doivent être apportés pour éduquer les individus, notamment en modifiant les programmes d'études et en discutant des pratiques autochtones et des améliorations à apporter à la formation des ingénieurs.

\stanceheading{La position des étudiants}
\begin{itemize}
 \item La représentation des autochtones dans le secteur de l'ingénierie est insuffisante si l'on considère uniquement les ingénieurs en âge de travailler (25-64 ans) ayant un niveau d'études égal ou supérieur au niveau du baccalauréat, seuls 0,93 \% d'entre eux s'identifient comme autochtones [10].
 \begin{itemize}
  \item En 2021, 49,2 \% des autochtones en âge de travailler avaient obtenu un certificat, un titre ou un diplôme postsecondaire, un taux inférieur à celui des non-autochtones (68 \%) [11].
  \item En 2021, les peuples autochtones représentaient 5,0 \% de la population totale du Canada, soit environ 1,8 million de personnes [12].
 \end{itemize}
 \item L'augmentation de la représentation des populations autochtones dans les établissements d'enseignement postsecondaire devrait être une priorité absolue avec l'aide du gouvernement pour l'enseignement postsecondaire.
 \item Une représentation accrue dans la profession d'ingénieur peut répondre aux appels à l'action de la Commission de vérité et réconciliation du Canada.
 \item Les connaissances et la sagesse autochtones doivent constituer un élément fondamental de la formation des ingénieurs, afin d'apporter une perspective, des connaissances traditionnelles et des pratiques durables.
 \item Les ingénieurs interagissent avec les communautés autochtones et ont un impact direct sur elles par le biais du développement, de l'entretien et de la modification des infrastructures et des projets d'ingénierie ; il faut donc trouver une compréhension interculturelle fondée sur le respect mutuel pour promouvoir une communication positive et contribuer à la réconciliation.
\end{itemize}

\stanceheading{Les recommandations aux partenaires, aux parties intéressées et aux autres entités}
\begin{itemize}
 \item La FCEG demande à Ingénieurs Canada d'entendre les voix des étudiants autochtones en génie afin d'améliorer la représentation des peuples autochtones.
 \item La FCEG demande à Ingénieurs Canada d'impliquer la FCEG dans le sous-réseau de la DIEEN (\textit{Decolonizing and Indigenizing Engineering Education Network}) afin d'accroître les connaissances de la FCEG en matière de ressources et d'enquêtes.
 \item La FCEG demande au BCAPG d'améliorer les attributs actuels des diplômés et de les élargir pour y inclure les pratiques autochtones locales et l'éducation afin d'inclure la communication avec les communautés autochtones ainsi que le savoir et la sagesse autochtones.
 \item La FCEG demande aux doyens des facultés d'ingénierie du Canada d'inclure l'éducation et les pratiques autochtones dans les plans stratégiques, afin de s'assurer que les établissements accordent une grande importance à ces pratiques lorsqu'ils prennent des décisions.
 \item La FCEG demande aux DDIC d'intégrer l'éducation et les pratiques autochtones dans leurs établissements par le biais d'activités, de conférences et d'événements, à l'instar d'autres établissements d'enseignement postsecondaire.
 \item La FCEG demande aux DDIC de promouvoir le recrutement d'autochtones dans le domaine de l'ingénierie, de fournir des ressources et de rendre l'ingénierie accessible aux peuples autochtones.
\end{itemize}

\stanceheading{Ce que fait la FCEG}
\begin{itemize}
 \item Les conférences de la FCEG mettent actuellement l'accent sur l'éducation et la pratique autochtones par le biais de conférenciers, d'ateliers et de diverses activités.
 \item La FCEG participe au projet iBET PhD qui soutient les professeurs autochtones et noirs dans les domaines de l'ingénierie et de l'informatique au Canada.
 \item La FCEG utilise l'enquête nationale pour créer un paysage démographique, y compris l'identification des peuples autochtones et la participation à l'ingénierie.
\end{itemize}

\stanceheading{Ce que la FCEG envisage de faire}
\begin{itemize}
 \item Établir des liens avec la communauté de pratique d'Ingénieurs Canada appelée \textit{Decolonizing and Indigenizing Engineering Education Network} (DIEEN) pour traiter de l'accès des autochtones à la formation en ingénierie ainsi que de la vérité et de la réconciliation dans le cadre de la formation en ingénierie [13].
 \item Recueillir de l'information auprès des DDIC des établissements qui ont mis sur pied des initiatives d'ingénierie postsecondaire pour promouvoir l'éducation autochtone afin d'élaborer un cadre de travail pour la promotion de l'éducation autochtone (conformément à la documentation d'Ingénieurs Canada) :
 \begin{itemize}
  \item Cours en ligne sur les peuples autochtones et les technosciences à l'Université de l'Alberta
  \item \textit{Indigenous Futures in Engineering} à l'Université Queen's
  \item Programme d'accès au génie de l'Université du Manitoba
  \item Programme d'accès à l'ingénierie pour les autochtones de l'Université de la Saskatchewan
 \end{itemize}
 \item Travailler avec d'autres organisations partenaires, telles que l'\textit{American Indian Science and Engineering Society} (AISES), afin de mieux comprendre la situation et d'établir des partenariats en vue de poursuivre les efforts de sensibilisation.
 \item Analyser les informations et les enquêtes actuelles recueillies par le \textit{DIEEN} et Ingénieurs Canada afin de créer un paramètre initial pour la collecte de données et l'analyse qualitative/quantitative en vue d'intégrer des questions sur l'éducation autochtone dans l'enquête nationale.
\end{itemize}

\newpage
